\section{Kaggle lab: ensemble regression on Ames housing}
\label{sec:kaggle-house-prices-ensembles}

\begin{kaggleproject}{Ames Housing --- random forest versus gradient boosting}
\kagglesource{https://www.kaggle.com/datasets/shashanknecrothapa/ames-housing-dataset}{Ames Housing Dataset}
\textbf{File.} \texttt{all\_datasets/ames\_housing\_dataset/AmesHousing.csv}\\
\textbf{Target.} house sale price.\\
\textbf{Question.} How do bagging and boosting compare when they receive the same mixed-type preprocessing and cross-validation folds?\\
\textbf{Before running.} Predict which family may be more robust with minimal tuning and which may gain more from careful hyperparameter control.

\begin{labmode}
Stage 3B: preserve the same folds and preprocessing across model families. The Ames comparison is development-only cross-validation; the wine transfer uses a development validation holdout. Neither is presented as a final untouched-test estimate. Design the fair comparison first; then use the reference scripts to verify it.
\end{labmode}

\lstinputlisting[style=mlpython]{all_experiments/lab12-house_prices_ensembles.py}

\begin{pitfall}
Do not choose the better model from training error. The comparison uses out-of-fold validation scores produced by the same cross-validation design.
\end{pitfall}
\end{kaggleproject}

\subsection{Transfer lab: the same ensemble ideas on wine measurements}

\begin{kaggleproject}{Red Wine Quality --- transfer without categorical preprocessing}
\kagglesource{https://www.kaggle.com/datasets/uciml/red-wine-quality-cortez-et-al-2009}{Red Wine Quality}
\textbf{File.} \texttt{all\_datasets/red\_wine\_quality\_dataset/winequality-red.csv}\\
\textbf{Target.} integer quality score, treated here as regression.\\
\textbf{Question.} When the schema changes to a compact all-numeric table, do the same ensemble concepts still apply even though the preprocessing becomes simpler?

\lstinputlisting[style=mlpython]{all_experiments/lab12-wine_quality_ensembles.py}

\begin{tryit}
Inspect the target distribution first. Then compare regression with an alternative classification framing such as ``quality at least 6.'' Explain why changing the target definition changes the scientific question, metrics, and error costs --- it is not merely another parameter setting.
\end{tryit}
\end{kaggleproject}

\begin{labdeliverable}
Submit one same-fold comparison of random forest and gradient boosting on Ames Housing, including mean and spread of validation error. Then include the Red Wine development-validation transfer result and a short memo explaining how changing from regression to a ``quality at least 6'' classification target changes the scientific question rather than merely the model setting.
\end{labdeliverable}
